\begin{recipedef}
  Name:                & \REC{metis_n_img_chopnod} \\
  Purpose:             & Basic science image data processing \\
%  Requirements:        & \REQ{METIS-XXXX} \\
%  Type:                & Calibration | Science | Technical \\
%  Templates:           & \TPL{...} \\
  Primary input:       & \RAW{N_IMAGE_SCI_RAW} or \RAW{N_IMAGE_STD_RAW} (1--N) \\
  Secondary inputs:    & \PROD{BADPIX_MAP_GEO} (optional) \\
                       & \hyperref[dataitem:gain_map_det]{\STATCALIB*{GAIN_MAP_GEO}} \\
                       & \hyperref[dataitem:linearity_det]{\PROD*{LINEARITY_GEO}} \\
                       & \PROD{MASTER_DARK_GEO} \\
                       & \PROD{MASTER_IMG_FLAT_LAMP_N} or \PROD{MASTER_IMG_FLAT_TWILIGHT_N} \\
                       & \EXTCALIB{PERSISTENCE_MAP} (optional) \\
  Matched keywords:    & \FITS{DET.DIT} \\
                       & \FITS{DET.NDIT} \\
                       & \FITS{DRS.FILTER} \\
  Parameters:          & \CODE{metis_n_img_chopnod.stacking.method}\newline Name of the method used to combine the input images (\texttt{add}, \texttt{average}, \texttt{median}), default \texttt{add} \\
  Algorithm:           & Remove crosstalk, correct non-linearity \\
                       & Analyse and optionally remove masked regions \\
                       & Subtract dark, divide by flat \\
                       & Remove blank sky pattern \\
  Output data:         & \PROD{N_SCI_BKG_SUBTRACTED} or \PROD{N_STD_BKG_SUBTRACTED} \\
%  Expected accuracies: & \\
  QC1 parameters:      & \QC{QC N IMG PEAK CNTS} \\
\end{recipedef}

